
\subsubsection{Research Questions}

\textbf{RQ1 (Existence)}: Can the methodology detect distinguishable performance patterns when simulated with differentiated signatures (Cohen's d > 1.$5\sigma )$?

\textbf{RQ2 (Execution Mechanism)}: Can the framework identify execution-dominance patterns in correctness dimension (top-15\% percentile rank)?

\textbf{RQ3 (Preference Mechanism)}: Can the framework identify balanced performance patterns across all dimensions (mean rank $\leq 30$\%)?

\subsubsection{Datasets and Tasks}

Models are evaluated on HumanEval+, an extension of HumanEval containing 164 hand-crafted Python programming tasks at the function level, each with extended test suites (80+ tests per task on average). HumanEval+ was selected for standardization (de facto standard for alignment method evaluation), extended test coverage (reduces variance), function-level scope (controls complexity confounds), and public availability.

\subsubsection{Model Selection}

The proof-of-concept validation used 3 models:

\begin{itemize}
\item \textbf{microsoft/phi-2} (2.7B parameters): Instruction-tuned model with execution-based alignment
\item \textbf{Salesforce/codegen-350M-mono} (350M parameters): Analyzed for preference-like patterns
\item \textbf{Salesforce/codegen-350M-nl} (350M parameters): Baseline model
\end{itemize}

This selection balances alignment method diversity with practical constraints. The small sample size (3 models total) limits statistical power. A critical confound: phi-2 (2.7B) is $8\times$ larger than codegen-350M, potentially conflating alignment method effects with capacity effects.

\subsubsection{Evaluation Protocol}

\#\#\#\# Code Generation

For each model m and task t:

\begin{enumerate}
\item Load task description and function signature
\item Generate k=10 code samples at temperature T=0.8, top-p p=0.95
\item Extract function implementation from each sample
\item Save samples
\end{enumerate}

The minimal validation used 10 tasks (reduced from 164 for proof-of-concept feasibility) with k=10 samples per task.

\#\#\#\# Correctness Measurement

For each generated sample:

\begin{enumerate}
\item Execute sample against HumanEval+ extended test suite
\item Record pass/fail for each test case
\item Compute per-sample correctness: fraction of tests passed
\item Aggregate: pass@k rate (fraction of samples passing all tests)
\end{enumerate}

\#\#\#\# Complexity Measurement

For each syntactically valid sample:

\begin{enumerate}
\item Compute cyclomatic complexity using radon (CC = E - N + 2P where E=edges, N=nodes, P=connected components)
\item Compute AST depth via ast module (maximum nesting level)
\item Average across valid samples
\end{enumerate}

\#\#\#\# Efficiency Measurement

For each sample that passes all tests:

\begin{enumerate}
\item Execute sample with cProfile to measure runtime (milliseconds)
\item Measure peak memory usage via tracemalloc (MB)
\item Average across passing samples
\end{enumerate}

Efficiency is measured only on correct implementations to avoid noise from crashes or infinite loops.

\subsubsection{Signature Detection Analysis}

Given performance signatures \textbf{s}$_{m}$ for each model:

\begin{enumerate}
\item Standardize features via StandardScaler (zero mean, unit variance)
\item Apply PCA to compute principal components
\item Run k-means with k=3 clusters, 10 random restarts
\item Compute alignment purity and Cohen's d
\end{enumerate}

Success criterion (H-E1): d > 1.$5\sigma$ and alignment purity > 0.7.

For mechanistic validation:

\begin{enumerate}
\item Compute percentile ranks for each model on each dimension
\item Test M1 (execution dominance): Check if execution models achieve rank $\leq 15$\% on correctness
\item Test M2 (preference balance): Check if preference models achieve mean rank $\leq 30$\% across all dimensions
\end{enumerate}

\subsubsection{Implementation Details}

\#\#\#\# Proof-of-Concept Validation

The primary proof-of-concept validation used simulated performance data to validate the methodology pipeline (clustering, PCA, effect size computation). Simulated data was designed with differentiated signatures (execution models high correctness, preference models balanced, baselines low overall) to test whether the analysis pipeline detects these patterns.

This approach validates that signature detection methodology works (the statistical framework can detect signatures when they exist) but limits claims about real-world alignment method behavior until validated with real model inference.

\#\#\#\# Minimal Real-Model Validation

A minimal real-model validation was conducted using 3 models on 10 HumanEval+ tasks with 10 samples per task (300 total code generations). This confirmed the framework can profile actual models and extract real performance signatures, though the limited scale prevents robust statistical conclusions.

\#\#\#\# Computational Resources

Full real-model inference would require approximately 6,000 generations (164 tasks $\times$ 30 samples $\times$ 4 models) on a single GPU, estimated at 2-4 hours total runtime. Proof-of-concept simulation and minimal validation enabled rapid methodology validation.
